% =============================================================================
% §3 Motivation  — Target: 1.4--1.6 pages. Every finding is (A) exists (B) hurts.
% 4.5/5 bar: negative results become design constraints, not embarrassment.
% 填写: 换 Fig.~3--5 和 Table~1 的实测；caption 里的定性结论不要改，除非 finding 变了。
% =============================================================================
\section{Why Naive CXL-SSD ANNS Fails}
\label{sec:motivation}

A unified VA does not make graph ANNS fast. This section records the
pathologies that a runtime must treat as first-class. Each finding answers
two questions: \emph{does the effect exist in the numbers?} and \emph{how
much does being wrong cost?} We later require every design module to name
the finding it answers (Table~\ref{tab:findings}).

\subsection{F1: A two-order miss cliff}
\label{sec:mot-cliff}

A full-precision load that misses the CXL-SSD window costs
tens of microseconds on the measured path (true-cold mean promote
24--72\,$\mu$s; Figure~\ref{fig:mot-cliff}). A hit that is already in
the 1\,GiB window is a host memcpy (sub-microsecond). The ratio is two
orders of magnitude. Graph search multiplies this tax by hops
$\times$ degree. Once a modest fraction of loads miss, wall-clock time is
spent in the NAND path, not in distance arithmetic. DiskANN already knew
this for block I/O; the new fact is that \emph{memory-semantic} access does
not remove the cliff --- it only changes how a miss is issued
(\S\ref{sec:mot-qd}).

\begin{figure}[t]
  \centering
  \phbox{0.98\columnwidth}{3.3cm}{%
    Fig.~3 placeholder --- three bars or a CDF:\\
    cold NAND / soft-cache$+$PTE-cold / PTE-warm.\\
    Optional stacked bar: fraction of a full search in each tier.}
  \caption{Miss cliff on flash-backed CXL memory. Soft-cache hits are
    \emph{not} PTE hits; reporting only \texttt{mincore}/PTE occupancy
    hides the middle step. \textbf{Fill-in:} per-vector latency and a
    whole-search breakdown on the same index.}
  \Description{Bar chart of vector-load latency at three residency tiers.}
  \label{fig:mot-cliff}
\end{figure}

\subsection{F2: Page admission mismatches the graph}
\label{sec:mot-page}

Neighbor IDs are not sequential in the address space. Admitting 16\,KB OS
pages (or CXL-SSD cachelines grouped as pages) therefore brings vectors the
search will not score. On a high-dimensional corpus where a vector is
already close to a page, page admission versus per-record admission
splits QPS by \num{$\sim$3.6$\times$} and NAND reads by
\num{$\sim$3.7$\times$} (Figure~\ref{fig:mot-pathology}a). The claim is
scoped: at small dimensions several records share a page and the gap
narrows. Default page-cache thinking is still the wrong default for this
workload.

\subsection{F3: Promotion volume without precision is harmful}
\label{sec:mot-pref}

Synchronous logical prefetch of 1-, 2-, or 8-hop neighborhoods leaves
hit rate almost unchanged, multiplies NAND reads, and drops QPS by
\num{$\sim$4--5$\times$} at the widest setting
(Figure~\ref{fig:mot-pathology}b). Precision (\emph{used$/$promoted})
falls as coverage grows; leftover pages evict useful ones from a
few-gigabyte window. \emph{Wrong promotion costs more than a missed
promotion.} The useful policy is a bounded, high-precision ensure of the
nodes the search is about to expand --- not ``prefetch harder.''

\subsection{F4: Search width is a stall knob}
\label{sec:mot-beam}

Raising \texttt{efSearch} / beam grows the unique IDs touched per query,
cools the window, and raises NAND reads faster than it raises recall
(Figure~\ref{fig:mot-pathology}c). Recall itself saturates once a hop
budget is spent. A shared hot set at a \emph{small} beam beats a cold
cache at a \emph{large} beam. A controller that chops beam on instantaneous
miss rate does not, by itself, win at iso-recall: it has no recall
constraint and no IO budget. Quality and stall must be controlled jointly.

\begin{figure*}[t]
  \centering
  \phbox{0.98\textwidth}{3.5cm}{%
    Fig.~4 placeholder --- three panels, same y-family (QPS or normalized):\\
    (a)~record vs.\ 16\,KB page admission;\\
    (b)~demand vs.\ 1-/2-/8-hop sync prefetch (precision annotated);\\
    (c)~efSearch / beam sweep: QPS, hit\%, recall@10, NAND reads.}
  \caption{Three default usages that amplify stall. (a)~Page admission.
    (b)~Blind graph prefetch. (c)~Widening search. Each panel must keep
    recall visible so raw QPS cannot be misread.
    \textbf{Fill-in:} one protocol (dataset, $L$, cache bytes) across
    panels if possible.}
  \Description{Three-panel figure of page admission, prefetch precision,
    and search-width coupling.}
  \label{fig:mot-pathology}
\end{figure*}

\subsection{F5: A single query cannot fill the miss pipeline}
\label{sec:mot-qd}

Graph search is pointer-chasing. Even an excellent single-query pipeline
leaves device-queue bubbles: the next hop is unknown until the current
page is scored, the beam ramps slowly, and the query drains at the end
(Figure~\ref{fig:mot-bubbles}). On NVMe, interleaving independent queries
raises occupancy; a \emph{continuous}, fetch-level scheduler that admits
work at read granularity can hold a target depth $D$.

On CXL-SSD the same idea is necessary and insufficient. A blocking
\texttt{memcpy} from a mapping is a synchronous NAND fault (effective
QD$\approx 1$ per thread). A depth-oriented scheduler that issues
\emph{speculative} fills then pays their latency on the critical path and
can lose to a narrower pipeline. Continuous batching pays off on this tier
only after memory-semantic fills are asynchronous
(Figure~\ref{fig:design-pipe}, \S\ref{sec:d4}).

\begin{figure}[t]
  \centering
  \phbox{0.98\columnwidth}{3.4cm}{%
    Fig.~5 placeholder --- in-flight depth vs.\ time:\\
    (a)~single-query pipeline; (b)~static batch; (c)~continuous fetch batching.\\
    (Existing TikZ in sections/fig\_bubbles.tex can replace this box.)}
  \caption{Pipeline bubbles. $y$ is in-flight device reads; $D$ is the
    depth needed to exercise NAND parallelism. A single query cannot hold
    $D$; a static batch drains at the barrier; continuous fetch-level
    admission targets $D$ at unchanged recall.
    \textbf{Fill-in:} swap in the measured schematic; add a CXL-SSD panel
    that shows blocking load/store stuck at QD$\approx 1$.}
  \Description{Three timelines of device queue occupancy under different
    schedulers.}
  \label{fig:mot-bubbles}
\end{figure}

\subsection{From findings to design constraints}
\label{sec:mot-map}

Table~\ref{tab:findings} is the contract for~\S\ref{sec:design}. We do not
treat CXL-SSD as uniform memory (C1). We place the graph and PQ off the
SSD window (C2). We promote few, precise records (C3) and share them
across queries (C4). We control beam under iso-recall or an IO budget
(C5), and we issue misses asynchronously at a target depth. Evaluation
must compare against a block-SSD path on the same corpus and recall
(C6).

\begin{table}[t]
  \caption{Findings become challenges become knobs.
    \textbf{Fill-in:} replace \ph\ with the headline measured delta
    (e.g., QPS $\times$, IO $\times$).}
  \label{tab:findings}
  \small
  \begin{tabular}{@{}p{0.07\columnwidth}p{0.28\columnwidth}p{0.26\columnwidth}p{0.27\columnwidth}@{}}
    \toprule
    ID & Finding (exists $+$ hurts) & Challenge & \sys{} knob \\
    \midrule
    F1 & Miss cliff \ph$\times$ & C1 visibility of tiers & counters, cost model \\
    F2 & Page $\neq$ record \ph$\times$ & C2 placement & record / hub admission \\
    F3 & Low-precision promote \ph$\times$ & C3 bounded promote & top-$M$, byte budget \\
    F4 & Beam cools the set \ph & C5 iso-recall & beam / early stop \\
    F5 & Bubbles; mmap QD$\approx 1$ & C4/C5 pipeline & async fill $+$ \texttt{cont} \\
    F7 & Shared hot set $>$ deep beam & C4 shared cache & cross-query pin \\
    \bottomrule
  \end{tabular}\\[2pt]
  {\footnotesize \phnote}
\end{table}
